%% ============================================================================
%% Fichier  : annexes/Ann_D_textes.tex
%% Titre    : Textes Juridiques et Normatifs
%% Version  : 1.0 -- novembre 2026 (INF3226, d'après ACSI)
%% ============================================================================
%% Extraits identiques à ceux des chapitres. Loi n° 2010/012 : texte
%% français consulté. COBAC R-2016/04 : relevé sur une base en ligne, à
%% collationner. ISSAI 400.35 : citée d'après le manuel de l'IDI (2021).
%% GUID 5100 § 6.10 : paraphrase vérifiée dans le texte français.
%% ============================================================================

\chapter{Textes Juridiques et Normatifs}
\label{ann:textes}

Cette annexe rassemble les textes cités dans le cours, avec leur degré de
vérification. Seuls les extraits lus dans leur version officielle sont
reproduits entre guillemets (Chapitre~\ref{chap:droit}).

\section*{Degré de vérification}

\begin{tableaucours}{L{4.6cm}L{3.6cm}L{4.0cm}}{%
    \tableauHeader{Texte} &
    \tableauHeader{Vérification} &
    \tableauHeader{Chapitres}
}
Loi n°~2010/012 (cybersécurité) & Texte officiel lu~; extraits cités &
\ref{chap:audit-auditologie}, \ref{chap:droit},
\ref{chap:audit-smsi}, \ref{chap:intrusion}, \ref{chap:audit-preuve} \\
Loi n°~2010/013 (communications électroniques) & Présentée~; articles à
vérifier & \ref{chap:droit} \\
Règlement COBAC R-2016/04 & Art.~1 et~2 relevés~; à collationner &
\ref{chap:controle-interne}, \ref{chap:droit} \\
AUSCGIE révisé (2014) & Présenté sans citation~; articles à vérifier &
\ref{chap:fraude}, \ref{chap:droit} \\
Acte uniforme sur le droit comptable (2017) & Présenté~; à vérifier &
\ref{chap:droit} \\
Loi n°~2018/012 (régime financier) & D'après la présentation du MINFI &
\ref{chap:droit} \\
Loi n°~2024/017 (données personnelles) & À vérifier au Journal officiel &
\ref{chap:droit}, \ref{chap:audit-conformite}, \ref{chap:audit-ia} \\
Code des marchés publics & Texte en vigueur à vérifier &
\ref{chap:audit-projets} \\
ISSAI 400 et 4000 (INTOSAI) & Cités d'après le manuel de l'IDI &
\ref{chap:normes} \\
GUID 5100 (INTOSAI) & Texte français lu & \ref{chap:normes},
\ref{chap:audit-continu} \\
\end{tableaucours}

\section*{Loi n° 2010/012 relative à la cybersécurité et à la cybercriminalité}

\begin{boxjuris}
\textbf{Article~4, point~9}~: \textit{« Audit de sécurité~: examen
méthodique des composantes et des acteurs de la sécurité, de la politique,
des mesures, des solutions, des procédures et des moyens mis en oeuvre par
une organisation, pour sécuriser son environnement, effectuer des contrôles
de conformité, des contrôles d'évaluation de l'adéquation des moyens
(organisationnels, techniques, humains, financiers) investis au regard des
risques encourus, d'optimisation, de rationalité et de performance. »}

\medskip
\textbf{Article~13, alinéa~1}~: \textit{« Sont soumis à un audit de sécurité
obligatoire, les réseaux de communications électroniques et les systèmes
d'information, des opérateurs, des autorités de certification et des
fournisseurs de services de communications électroniques. »}

\medskip
\textbf{Article~14}~: \textit{« Le personnel de l'Agence et les experts
commis en vue d'accomplir des opérations d'audit sont astreints au secret
professionnel. »}

\medskip
\textbf{Article~26, alinéa~1}~: \textit{« Les exploitants des systèmes
d'information prennent toutes les mesures techniques et administratives afin
de garantir la sécurité des services offerts. A cet effet, ils se dotent de
systèmes normalisés leur permettant d'identifier, d'évaluer, de traiter et
de gérer continûment les risques liés à la sécurité des systèmes
d'information dans le cadre des services offerts directement ou
indirectement. »}

\medskip
\textbf{Article~29, alinéa~1}~: \textit{« Les exploitants des systèmes
d'information ont l'obligation de conserver les données de connexion et de
trafic de leurs systèmes d'information pendant une période de dix (10)
ans. »}

\medskip
\textbf{Article~32, alinéas~1 à~3}~: \textit{« (1) Les réseaux de
communications électroniques et les systèmes d'information sont soumis à un
régime d'audit de sécurité obligatoire et périodique de leurs systèmes de
sécurité par l'Agence. (2) L'audit de sécurité et les mesures d'impact de
gravité sont effectués chaque année ou lorsque les circonstances
l'exigent. (3) Les rapports d'audit sont confidentiels et adressés au
Ministre chargé des Télécommunications. »}

\medskip
\textbf{Article~61, alinéa~1}~: \textit{« Sont punis d'un emprisonnement de
trois (03) mois à trois (03) ans et d'une amende de 20.000 (vingt mille) à
100.000 (cent mille) F CFA, les personnels de l'Agence et les experts des
personnes morales chargés des audits qui révèlent sans autorisation, des
informations confidentielles dont ils ont eu connaissance à l'occasion d'un
audit de sécurité. »}

\medskip
\textbf{Article~69}~: \textit{« Est puni d'un emprisonnement de cinq (05) à
dix (10) ans et d'une amende de 10.000.000 (dix millions) à 100.000.000
(cent millions) F CFA ou de l'une de ces deux peines seulement, celui qui
accède sans droit, et en violation des mesures de sécurité, à l'ensemble ou
à une partie d'un réseau de communications électroniques, d'un système
d'information ou d'un équipement terminal, afin d'obtenir des informations
ou des données, en relation avec un système d'information connecté à un
autre système d'information. »}
\end{boxjuris}

\begin{boxjuris}
\textbf{Dispositions citées dans les Chapitres~\ref{chap:intrusion}
et~\ref{chap:audit-preuve}} (texte officiel)

\medskip
\textbf{Article~65, alinéa~2}~: \textit{« Est puni des peines prévues à
l'alinéa 1 ci-dessus, tout accès non autorisé, à l'ensemble ou à une
partie d'un réseau de communications électroniques ou d'un système
d'information ou d'un équipement terminal. »}

\medskip
\textbf{Article~70}~: \textit{« Est puni d'une amende de 1.000.000 (un
million) à 5.000.000 (cinq millions) F CFA, celui qui provoque par
saturation, l'attaque d'une ressource de réseau de communications
électroniques ou d'un système d'information dans le but de l'effondrer en
empêchant la réalisation des services attendus. »}

\medskip
\textbf{Article~86, alinéa~1}~: \textit{« Est puni des peines prévues
l'article 71 ci-dessus, celui qui importe, détient, offre, cède, vend ou
met à disposition, sous quelle que forme que ce soit, un programme
informatique, un mot de passe, un code d'accès ou toutes données
informatiques similaires conçus et ou spécialement adaptés, pour permettre
d'accéder, à tout ou partie d'un réseau de communications électroniques ou
d'un système d'information. »}

\medskip
\textbf{Article~52, alinéa~3}~: \textit{« Les Officiers de Police Judiciaire
et les agents habilités de l'Agence peuvent, lors des investigations,
accéder aux moyens de transport, à tout local à usage professionnel, à
l'exclusion des domiciles privés, en vue de rechercher, de constater les
infractions, de demander la communication de tous les documents
professionnels et en prendre copie, recueillir, sur convocation ou sur
place, les renseignements et justifications. »}

\medskip
\textbf{Article~53, alinéa~1}~: \textit{« Les perquisitions en matière de
cybercriminalité sont susceptibles de porter sur des données qui peuvent
être des supports physiques ou des copies réalisées en présence des
personnes qui assistent à la perquisition. »}
\end{boxjuris}

Sont seulement résumés dans le cours~: l'article~7, alinéa~2 (missions de
l'ANTIC), l'article~30 (révision des systèmes de sécurité), l'article~61,
alinéa~3 (obstruction aux audits), les articles~54 et~55 (procédure,
réquisition pour le déchiffrement) et les articles~66 à~68 et~71 à~72
(atteintes aux systèmes et aux données).

\section*{Règlement COBAC R-2016/04}

\begin{boxjuris}
\textbf{Articles~1 et~2} (texte relevé sur une base réglementaire en ligne,
à collationner avec la version du SG-COBAC)~: article~1~: \textit{« Le
présent règlement fixe les exigences minimales relatives au système de
contrôle interne que les établissements assujettis doivent mettre en
œuvre. »}~; article~2, risque opérationnel~: \textit{« risque de pertes
résultant de l'inadéquation ou de la défaillance des processus internes, des
personnes et des systèmes, ou d'évènements externes. »}
\end{boxjuris}

\section*{Normes de l'audit public}

\begin{itemize}
    \item \textbf{ISSAI~400.35}, citée par le manuel de l'IDI~: l'audit de
          conformité \textit{« repose sur une relation tripartite dans
          laquelle l'auditeur vise à obtenir des éléments probants suffisants
          et appropriés afin de formuler une conclusion destinée à renforcer
          le degré de confiance des utilisateurs prévus, autres que la partie
          responsable, quant à la mesure ou à l'évaluation d'un sujet par
          rapport à des critères »}~\cite{idi2021conformite}.
    \item \textbf{ISSAI~4000.14} (résumé d'après le même manuel)~: on ne se
          réfère pas à la norme sans signaler les exigences non suivies et
          leurs conséquences.
    \item \textbf{GUID~5100, §~6.10} (résumé)~: usages des techniques
          d'audit assistées par ordinateur, et leurs avantages~: grands
          volumes, répétabilité, documentation horodatée~\cite{intosai2019guid5100}.
\end{itemize}
